\label{sec:background}

\subsection{LODES Workplace Area Characteristics}

The Longitudinal Employer-Household Dynamics (LEHD) program at the U.S.\
Census Bureau compiles employment records from state Unemployment Insurance
(UI) wage files and matches them to address-level employer geocodes
\citep{abowd2009lodes}.
The result is a block-level employment dataset with near-complete private-sector
coverage (UI-covered employment, approximately 95\% of all wage and salary jobs).

The Workplace Area Characteristics (WAC) file is one of three LODES output
tables.
It answers the question: ``For each census block, how many jobs are located
there, broken down by industry?''
This is the workplace perspective---where people work, not where they live
(the Origin Area Characteristics, OAC, answers the residential question).

\paragraph{WAC structure.}
Each row in the WAC file represents one census block.
The primary columns are:
\begin{itemize}
  \item \texttt{w\_geocode} (15-character block GEOID)
  \item \texttt{C000} (total jobs, all types)
  \item \texttt{CNS01}--\texttt{CNS20} (jobs by NAICS supersector)
\end{itemize}

The 20 NAICS sectors cover the full private economy from agriculture
(CNS01) to public administration (CNS20).
Table~\ref{tab:cns-sectors} lists the four sectors used for the commercial
intensity and industrial fraction signals.

\begin{table}[ht]
\centering
\caption{LODES WAC NAICS sector columns used for economic character computation.}
\label{tab:cns-sectors}
\small
\begin{tabular}{lll}
\toprule
Column & NAICS description & Signal \\
\midrule
CNS07 & Retail Trade & Commercial intensity \\
CNS09 & Information & Commercial intensity \\
CNS10 & Finance and Insurance & Commercial intensity \\
CNS11 & Real Estate and Rental and Leasing & Commercial intensity \\
\midrule
CNS01 & Agriculture, Forestry, Fishing, and Hunting & Industrial fraction \\
CNS02 & Mining, Quarrying, and Oil and Gas Extraction & Industrial fraction \\
CNS05 & Manufacturing & Industrial fraction \\
CNS08 & Transportation and Warehousing & Industrial fraction \\
\bottomrule
\end{tabular}
\end{table}

\paragraph{Block-to-tract aggregation.}
Census block GEOIDs are 15 characters: state (2) + county (3) + census tract
(6) + block (4).
Aggregating to census tracts requires summing all block-level counts within
each unique 11-character prefix.
After aggregation, each tract has a total job count \texttt{C000} and sector
subtotals \texttt{CNS01}--\texttt{CNS20}.

\paragraph{Geographic scope and availability.}
LODES WAC is available for all 50 states (plus DC and Puerto Rico) from 2002
through 2021 in LODES version~8 \citep{lodes2022}.
Download format: gzip-compressed CSV, one file per state per year.
The 2020 vintage is used in this paper.
Tract counts after aggregation from the 2020 run:
North Carolina 2,655; Wisconsin 1,527; Texas 6,877.

\paragraph{Tracts with zero jobs.}
Purely residential tracts with no employer establishments appear in the
LODES geographic coverage but have no WAC records.
These tracts receive a zero character vector $(0, 0, 0)$ by definition;
the cosine similarity for two zero-vector tracts is set to 1.0
(both residential, keep together).
This is the most common case in rural and low-density suburban areas.

\subsection{The Three Economic Character Signals}

\paragraph{Commercial intensity (CI).}
Measures the fraction of jobs in the service economy sectors that characterize
commercial districts: retail trade (CNS07), information technology and media
(CNS09), financial services (CNS10), and real estate (CNS11).
High CI values ($>$0.4) are found in downtown retail cores, suburban
strip-mall corridors, financial district nodes, and mixed-use commercial zones.
Low CI values ($<$0.05) are found in residential tracts, agricultural areas,
and industrial parks.

\paragraph{Industrial fraction (IF).}
Measures the fraction of jobs in goods-producing and logistics sectors:
agriculture and forestry (CNS01), extractive industries (CNS02),
manufacturing (CNS05), and transportation and warehousing (CNS08).
High IF values ($>$0.5) are found in factory towns, industrial parks,
port and logistics zones, and agricultural processing areas.
These are the zones where a significant fraction of the working population is
employed in production, extraction, or movement of physical goods.

\paragraph{Jobs per resident (JPR).}
Measures the total job count relative to tract population.
Values below 0.3 indicate bedroom suburbs where most residents commute out.
Values between 1.0 and 3.0 indicate employment centers that attract net
in-commuting.
Values above 5.0 indicate major employment concentrations: university research
parks, hospital campuses, data centers, airport zones.
JPR is capped at 10.0 to prevent numerical instability in tracts with
very high job density and very low population.

\paragraph{Why these three signals?}
The three signals together cover the primary economic character dimensions
that divide land into recognizably distinct zones:
the division between residential and workplace land (JPR),
the division between service economy and goods economy (CI vs.\ IF),
and the internal division within the service economy between
consumer-facing retail and professional services.
Together they produce meaningful separation between the four most common
tract economic types: bedroom suburb, commercial corridor, industrial park,
and mixed employment center.

\subsection{Geographic Coverage Adequacy}

LODES data has near-complete coverage of private-sector employment.
The primary coverage gaps are:
\begin{itemize}
  \item Federal civilian employees (excluded from UI wage records).
    Military bases and federal office complexes are undercounted.
  \item Self-employed workers (not covered by UI).
    Agricultural tracts with many self-employed farmers may have
    low apparent job counts.
  \item Tracts that cross state boundaries.
    LODES is processed per state; cross-state GEOID issues are rare but exist.
\end{itemize}
For the purpose of edge weight computation, these gaps are handled gracefully:
tracts with missing or zero WAC records receive zero character vectors and
are treated as residential (high mutual similarity with other zero-vector tracts).
